% GENERATED FILE. Edit canonical lessons, then run npm run generate:books.
% canonical-source-hash: fnv1a64:1be52844390310ed
% canonical-lessons: MR-C363-satta, MR-C363-bhrastacar, MR-C363-pardarsakta, MR-C363-svartha, MR-C363-cithavne

\chapter{Power and Its Abuse}
\label{ch:mr-power-and-its-abuse}
\hlchaptermodality{363}

\begin{chapteropening}
\textbf{By the end of this chapter:} \emph{I can talk about power, corruption and self-interest.}

Complete the last lesson of chapter 363: I can talk about power, corruption and self-interest.
\end{chapteropening}

\section[\texorpdfstring{\mr{सत्ता} (sattā)}{सत्ता (sattā)}]{\mr{सत्ता} (sattā) — power (political)}

\label{lesson:MR-C363-satta}



\begin{quote}
Before the new one: say the Marathi for to protest, then the Marathi for to boycott.
\end{quote}

\subsection*{You'll want to know: \mr{सत्ता}}
\textbf{\mr{सत्ता}} — \emph{sattā} — \textquotedblleft{}power (political)\textquotedblright{}.

\begin{quote}
\mr{सत्ता} \mr{हातात} \mr{आली} \mr{की} \mr{माणूस} \mr{बदलतो}.
\end{quote}

\begin{quote}
\emph{sattā hātāt ālī kī māṇūs badalto.}
\end{quote}

\begin{quote}
When power comes into his hands, a man changes.
\end{quote}

Feminine.

\mr{सत्तेवर} \mr{येणे}, to come to power.

\begin{cousinweb}
\textbf{\mr{सत्ता}} (\emph{sattā}) is a tatsama, taken from Sanskrit \emph{sattā}, \textquotedblleft{}existence, being\textquotedblright{} (Monier-Williams). Marathi uses it for the power to rule.
\end{cousinweb}

\begin{grammarlens}[title={What you've built}]
One more word for giving reasons and opinions.
\end{grammarlens}

\subsection*{Your turn}
\begin{itemize}
\raggedright
  \item \emph{Say it:} \emph{sattā}
  \item \emph{Say it:} \emph{sattā}, once more
  \item \emph{Say it:} \emph{sattā hātāt ālī kī māṇūs badalto.} — When power comes into his hands, a man changes.
  \item \emph{Recall:} say the Marathi for to protest, then the Marathi for to boycott, then say \emph{sattā} again
\end{itemize}

\begin{tcolorbox}[breakable,colback=teal!4,colframe=teal!35!black,title={Before you move on}]
What does \mr{सत्ता} mean? (\textquotedblleft{}Power (political)\textquotedblright{}.) Say it once more.
\end{tcolorbox}
\section[\texorpdfstring{\mr{भ्रष्टाचार} (bhraṣṭācār)}{भ्रष्टाचार (bhraṣṭācār)}]{\mr{भ्रष्टाचार} (bhraṣṭācār) — corruption}

\label{lesson:MR-C363-bhrastacar}



\begin{quote}
Before the new one: say the Marathi for to boycott, then the Marathi for power.
\end{quote}

\subsection*{You'll want to know: \mr{भ्रष्टाचार}}
\textbf{\mr{भ्रष्टाचार}} — \emph{bhraṣṭācār} — \textquotedblleft{}corruption\textquotedblright{}.

\begin{quote}
\mr{भ्रष्टाचारामुळे} \mr{रस्त्यांचं} \mr{काम} \mr{खराब} \mr{झालं}.
\end{quote}

\begin{quote}
\emph{bhraṣṭācārāmuḷe rastyāñcaṁ kām kharāb jhālaṁ.}
\end{quote}

\begin{quote}
Because of corruption the road work was poor.
\end{quote}

Masculine.

\mr{भ्रष्ट}, corrupt, + \mr{आचार}, conduct.

\begin{cousinweb}
\textbf{\mr{भ्रष्टाचार}} (\emph{bhraṣṭācār}) is a learned compound built from Sanskrit \emph{bhraṣṭa}, \textquotedblleft{}fallen, ruined\textquotedblright{}, and \emph{ācāra}, \textquotedblleft{}conduct\textquotedblright{} (Monier-Williams).
\end{cousinweb}

\begin{grammarlens}[title={What you've built}]
One more word for giving reasons and opinions.
\end{grammarlens}

\subsection*{Your turn}
\begin{itemize}
\raggedright
  \item \emph{Say it:} \emph{bhraṣṭācār}
  \item \emph{Say it:} \emph{bhraṣṭācār}, once more
  \item \emph{Say it:} \emph{bhraṣṭācārāmuḷe rastyāñcaṁ kām kharāb jhālaṁ.} — Because of corruption the road work was poor.
  \item \emph{Recall:} say the Marathi for to boycott, then the Marathi for power, then say \emph{bhraṣṭācār} again
\end{itemize}

\begin{tcolorbox}[breakable,colback=teal!4,colframe=teal!35!black,title={Before you move on}]
What does \mr{भ्रष्टाचार} mean? (\textquotedblleft{}Corruption\textquotedblright{}.) Say it once more.
\end{tcolorbox}
\section[\texorpdfstring{\mr{पारदर्शकता} (pārdarśaktā)}{पारदर्शकता (pārdarśaktā)}]{\mr{पारदर्शकता} (pārdarśaktā) — transparency}

\label{lesson:MR-C363-pardarsakta}



\begin{quote}
Before the new one: say the Marathi for power, then the Marathi for corruption.
\end{quote}

\subsection*{You'll want to know: \mr{पारदर्शकता}}
\textbf{\mr{पारदर्शकता}} — \emph{pārdarśaktā} — \textquotedblleft{}transparency\textquotedblright{}.

\begin{quote}
\mr{सरकारी} \mr{कामात} \mr{पारदर्शकता} \mr{हवी}.
\end{quote}

\begin{quote}
\emph{sarkārī kāmāt pārdarśaktā havī.}
\end{quote}

\begin{quote}
Government work needs transparency.
\end{quote}

Feminine.

\mr{पारदर्शक}, transparent, + -\mr{ता}.

\begin{cousinweb}
\textbf{\mr{पारदर्शकता}} (\emph{pārdarśaktā}) is a learned noun built from Sanskrit \emph{pāradarśaka}, \textquotedblleft{}showing the opposite shore\textquotedblright{} (Monier-Williams), with \emph{-tā}, \textquotedblleft{}-ness\textquotedblright{}: what lets one see through.
\end{cousinweb}

\begin{grammarlens}[title={What you've built}]
One more word for giving reasons and opinions.
\end{grammarlens}

\subsection*{Your turn}
\begin{itemize}
\raggedright
  \item \emph{Say it:} \emph{pārdarśaktā}
  \item \emph{Say it:} \emph{pārdarśaktā}, once more
  \item \emph{Say it:} \emph{sarkārī kāmāt pārdarśaktā havī.} — Government work needs transparency.
  \item \emph{Recall:} say the Marathi for power, then the Marathi for corruption, then say \emph{pārdarśaktā} again
\end{itemize}

\begin{tcolorbox}[breakable,colback=teal!4,colframe=teal!35!black,title={Before you move on}]
What does \mr{पारदर्शकता} mean? (\textquotedblleft{}Transparency\textquotedblright{}.) Say it once more.
\end{tcolorbox}
\section[\texorpdfstring{\mr{स्वार्थ} (svārtha)}{स्वार्थ (svārtha)}]{\mr{स्वार्थ} (svārtha) — self-interest}

\label{lesson:MR-C363-svartha}



\begin{quote}
Before the new one: say the Marathi for corruption, then the Marathi for transparency.
\end{quote}

\subsection*{You'll want to know: \mr{स्वार्थ}}
\textbf{\mr{स्वार्थ}} — \emph{svārtha} — \textquotedblleft{}self-interest\textquotedblright{}.

\begin{quote}
\mr{तो} \mr{फक्त} \mr{स्वतःचा} \mr{स्वार्थ} \mr{पाहतो}.
\end{quote}

\begin{quote}
\emph{to phakta svataḥcā svārtha pāhto.}
\end{quote}

\begin{quote}
He only looks after his own interest.
\end{quote}

Masculine.

\mr{स्वार्थी}, selfish; \mr{निःस्वार्थ}, selfless.

\begin{cousinweb}
\textbf{\mr{स्वार्थ}} (\emph{svārtha}) is a tatsama, taken from Sanskrit \emph{svārtha}, \textquotedblleft{}one's own aim, self-interest\textquotedblright{} (Monier-Williams).
\end{cousinweb}

\begin{grammarlens}[title={What you've built}]
One more word for giving reasons and opinions.
\end{grammarlens}

\subsection*{Your turn}
\begin{itemize}
\raggedright
  \item \emph{Say it:} \emph{svārtha}
  \item \emph{Say it:} \emph{svārtha}, once more
  \item \emph{Say it:} \emph{to phakta svataḥcā svārtha pāhto.} — He only looks after his own interest.
  \item \emph{Recall:} say the Marathi for corruption, then the Marathi for transparency, then say \emph{svārtha} again
\end{itemize}

\begin{tcolorbox}[breakable,colback=teal!4,colframe=teal!35!black,title={Before you move on}]
What does \mr{स्वार्थ} mean? (\textquotedblleft{}Self-interest\textquotedblright{}.) Say it once more.
\end{tcolorbox}
\section[\texorpdfstring{\mr{चिथावणे} (cithāvṇe)}{चिथावणे (cithāvṇe)}]{\mr{चिथावणे} (cithāvṇe) — to incite; to provoke}

\label{lesson:MR-C363-cithavne}



\begin{quote}
Before the new one: say the Marathi for transparency, then the Marathi for self-interest.
\end{quote}

\subsection*{You'll want to know: \mr{चिथावणे}}
\textbf{\mr{चिथावणे}} — \emph{cithāvṇe} — \textquotedblleft{}to incite; to provoke\textquotedblright{}.

\begin{quote}
\mr{कोणीतरी} \mr{गर्दीला} \mr{चिथावलं}.
\end{quote}

\begin{quote}
\emph{koṇītrī gardīlā cithāvlaṁ.}
\end{quote}

\begin{quote}
Someone incited the crowd.
\end{quote}

Past, with -\mr{ने} (he …): \mr{त्याने} \mr{चिथावलं}, \emph{tyāne cithāvlaṁ}.

The one incited takes -\mr{ला}; \mr{चिथावणी}, incitement.

\begin{cousinweb}
\textbf{\mr{चिथावणे}} (\emph{cithāvṇe}) is built on an older verb, Marathi \emph{cithṇẽ}, \textquotedblleft{}to take offence\textquotedblright{}, which CDIAL lists under \emph{*citth-}, \textquotedblleft{}pound, tear\textquotedblright{} (CDIAL 4802): to make someone take offence.
\end{cousinweb}

\begin{grammarlens}[title={What you've built}]
One more word for giving reasons and opinions.
\end{grammarlens}

\subsection*{Your turn}
\begin{itemize}
\raggedright
  \item \emph{Say it:} \emph{cithāvṇe}
  \item \emph{Say it:} \emph{cithāvṇe}, once more
  \item \emph{Say it:} \emph{koṇītrī gardīlā cithāvlaṁ.} — Someone incited the crowd.
  \item \emph{Recall:} say the Marathi for transparency, then the Marathi for self-interest, then say \emph{cithāvṇe} again
\end{itemize}

\begin{tcolorbox}[breakable,colback=teal!4,colframe=teal!35!black,title={Before you move on}]
What does \mr{चिथावणे} mean? (\textquotedblleft{}To incite; to provoke\textquotedblright{}.) Say it once more.
\end{tcolorbox}
